\documentclass[11pt]{article}
\usepackage{amsmath,amssymb}
\usepackage{bm}
\usepackage{geometry}
\geometry{margin=1in}

\title{Problem Formulation: Event-Driven Electric Truck Routing and Charging}
\author{}
\date{}

\begin{document}
\maketitle

\section{Routing Problem and Sources of Uncertainty}
We operate a fleet $\mathcal{T}=\{1,\dots,T\}$ of electric trucks over a road network $G=(\mathcal{V},\mathcal{E})$ with charging nodes $\mathcal{C}\subseteq\mathcal{V}$ and delivery nodes $\mathcal{D}\subseteq\mathcal{V}$. Each truck $i$ starts from a depot $d_{i0}$ and must serve a sequence of deliveries $\mathcal{D}_i=\{d_{i1},\dots,d_{iK}\}$ (sequentially or in flexible order, per the configuration). Battery capacity is $B$ (default $400$ kWh) with initial charge $b_i^{0}$ (``full'' or sampled). Travel is driven by base speeds $v$ (default $40$ km/h) and pre-computed shortest-path energy/time dictionaries $e_{uv}$, $\tau_{uv}$ for $(u,v)\in\mathcal{E}$.

\paragraph{Stochastic elements (from \texttt{config.yaml}).}
\begin{itemize}
    \item \textbf{Traffic and energy}: Travel time $\tilde{\tau}_{uv}$ is drawn from $\mathcal{N}(\tau_{uv},(\sigma_{\tau}\tau_{uv})^2)$ with rush-hour multiplier; energy consumption is $\tilde{e}_{uv}=e_{uv}\,\xi_{uv}$ where $\xi_{uv}$ is bounded in $[\underline{\xi},\overline{\xi}]$ (defaults $0.9$--$1.2$) and scales with traffic when \texttt{enable\_energy\_uncertainty} is on.
    \item \textbf{Unloading}: Service time $\tilde{u}_{d}$ at delivery $d$ follows $\mathcal{N}(u_0,(s u_0)^2)$ with business-hour variance multiplier and caps (\texttt{base\_unloading\_time}, \texttt{std\_dev\_factor}, \texttt{max\_std\_dev\_hours}).
    \item \textbf{Charging}: Chargers have type-dependent rates $r_c$ and efficiencies $\eta_c$; when \texttt{use\_realistic\_curve} is true, a CCCV taper applies after SOC $\geq \texttt{taper\_start\_soc}$, with minimum power \texttt{taper\_power\_min}. Available discrete durations are $h\in\mathcal{H}=\{1,\dots,10\}$ hours.
\end{itemize}

\paragraph{Environment dynamics.} The simulator is event-driven with a global clock and priority queue of events. Only the active truck (the one whose event is processed) is controlled. Event types are \textsc{TRUCK\_READY} (decision point) and \textsc{TRUCK\_ROUTING} (arrival). Truck states are \texttt{ready}, \texttt{routing}, \texttt{waiting\_to\_charge}, \texttt{charging}, \texttt{unloading}, \texttt{complete}, and \texttt{failed}. Charger ports enforce FCFS gating; trucks wait in queues and are woken when capacity frees. Episodes truncate at horizon $T_{\max}$ hours or step cap.

\section{Simplified Per-Truck Gurobi Model}
The baseline planner optimizes one truck at a time, ignoring charger contention and using nominal travel times. Let the remaining delivery sequence for truck $i$ be $(d_0,\dots,d_K)$ with $d_0$ the current node. For each leg $k=0,\dots,K-1$, options are to go directly or detour via one charger $c\in\mathcal{C}$.

\subsection*{Parameters}
\begin{itemize}
    \item $B$: battery capacity; $b^{\text{init}}$: current energy.
    \item $\alpha$: energy safety multiplier (=$\overline{\xi}$ if energy uncertainty enabled, else $1$).
    \item $e_{uv}$, $\tau_{uv}$: deterministic energy and time for $(u,v)$; $r_c,\eta_c$: charge rate and efficiency for charger $c$.
    \item $\mathcal{H}$: allowable charge durations (upper bound for continuous variables that are later rounded up).
    \item $M$: big-$M$ constant derived from worst-case energy and maximal charge addition.
\end{itemize}

\subsection*{Decision variables}
\begin{align*}
&b_k &&\text{battery before serving leg }k,\; k=0,\dots,K,\\
&x_k &&\in\{0,1\}\text{ direct-travel indicator for leg }k,\\
&y_{k,c} &&\in\{0,1\}\text{ use-charger indicator (at most one) for leg }k,\\
&t_{k,c} &&\ge 0\text{ continuous charge time at charger }c\text{ during leg }k,\\
&t_0 &&\ge 0\text{ optional initial charge time if }d_0\in\mathcal{C}.
\end{align*}

\subsection*{Objective}
\[
\min \; t_0 + \sum_{k=0}^{K-1}\Big( x_k \tau_{d_k,d_{k+1}} + \sum_{c\in\mathcal{C}} y_{k,c}(\tau_{d_k,c}+\tau_{c,d_{k+1}}) + \sum_{c\in\mathcal{C}} t_{k,c}\Big).
\]
This minimizes nominal travel plus charging time; solver tolerates a $10\%$ MIP gap and a $120$ s limit.

\subsection*{Constraints}
\begin{align*}
&b_0 =
    \begin{cases}
        b^{\text{init}} + r_c\eta_c t_0, & d_0=c\in\mathcal{C},\\
        b^{\text{init}}, & \text{otherwise},
    \end{cases}
    &&\text{(initial battery)}\\
&x_k + \sum_{c\in\mathcal{C}} y_{k,c} = 1, &&\forall k \quad\text{(one routing choice per leg)}\\
&b_k - \alpha e_{d_k,d_{k+1}} \ge 0 - M(1-x_k), &&\forall k \quad\text{(energy for direct leg)}\\
&b_{k+1} = b_k - \alpha e_{d_k,d_{k+1}} + M(1-x_k), &&\forall k \quad\text{(battery update if direct)}\\
&b_k - \alpha e_{d_k,c} \ge 0 - M(1-y_{k,c}), &&\forall k,c \quad\text{(reach charger)}\\
&b_{k+1} = b_k - \alpha e_{d_k,c} - \alpha e_{c,d_{k+1}} + \eta_c r_c t_{k,c} + M(1-y_{k,c}), &&\forall k,c \quad\text{(battery via charger)}\\
&0 \le t_{k,c} \le \max \mathcal{H}\, y_{k,c}, &&\forall k,c \quad\text{(charge only if chosen)}\\
&\underline{b} \le b_k \le B, &&\forall k \quad\text{(bounds, with } \underline{b}=0.1B \text{ buffer)}.
\end{align*}
If the truck starts outside any feasible leg (no direct path and no reachable charger), the model is infeasible and the policy falls back to an emergency plan (navigate to nearest reachable charger and over-charge to $95\%$). Continuous $t_{k,c}$ and $t_0$ are rounded up to the nearest duration in $\mathcal{H}$ (with a $10\%$ buffer) when actions are emitted to the environment.

\section{Markov Decision Process}
The single-agent controller observes a partially aggregated multi-truck system where only the active truck can act at a decision epoch.

\subsection*{MDP definition}
\begin{itemize}
    \item \textbf{State space $\mathcal{S}$}: $s=(g,i_{\text{act}},\{(n_j,b_j,\sigma_j,R_j)\}_{j\in\mathcal{T}},\chi)$, where $g$ is global clock, $i_{\text{act}}$ the active truck index, $n_j$ current node, $b_j$ battery, $\sigma_j\in\{\texttt{ready},\texttt{routing},\dots\}$ the truck mode, $R_j$ remaining delivery sequence (or set if flexible), and $\chi$ encodes charger occupancy/waitlists and pending events. The implemented observation vector mirrors this structure with normalized features for all trucks, deliveries, chargers, and global time.
    \item \textbf{Action space $\mathcal{A}(s)$}: see below.
    \item \textbf{Transition kernel $\mathcal{P}$}: Event-driven. If $a$ is a navigation to $v$, travel completes after random $\tilde{\tau}$ with battery update $b' = b - \tilde{e}$; if $a$ is charging for $h$, $b'$ follows the (possibly tapered) charging curve and completes after $h'$ returned by the curve model. Charger-gating transitions trucks to \texttt{waiting\_to\_charge} with queue-based wake-up. Unloading draws $\tilde{u}$ and holds the truck in \texttt{unloading} until completion. The next decision epoch occurs at the earliest \textsc{TRUCK\_READY} event.
    \item \textbf{Reward $r(s,a,s')$}: Time penalties $-\lambda_{\text{time}}\Delta t$ (with $\lambda_{\text{time}}=1$ by default) for travel, charging, waiting, and unloading; bonus $+R_{\text{del}}$ upon completing a delivery; terminal penalty $R_{\text{fail}}$ on failure (battery depletion or infeasible choice).
\end{itemize}

\subsection*{Feasible action set (key contribution)}
Actions are discrete, indexed as:
\begin{align*}
&\text{Navigation to charger }c\in\mathcal{C},\\
&\text{Navigation to delivery }d\in R_i \text{ (one action if sequential, one per remaining delivery if flexible)},\\
&\text{Charge at current node for }h\in\mathcal{H}.
\end{align*}
Feasibility masking enforces:
\begin{itemize}
    \item \textbf{Energy headroom}: Navigation to $v$ allowed only if $\alpha\,e_{n_i,v} < b_i$ and, for deliveries, there exists a follow-up charger reachable with remaining energy (unless the truck must leave a charger immediately).
    \item \textbf{Location constraints}: Charging actions valid only when $n_i\in\mathcal{C}$; navigation to current location is disallowed.
    \item \textbf{Charger gating}: If the truck just finished charging and must vacate the port, charging actions are blocked; otherwise, charging durations must yield $b_i + \Delta b(h) \ge$ energy required to depart safely.
    \item \textbf{Terminal states}: No actions are feasible for \texttt{complete}/\texttt{failed} trucks.
\end{itemize}
Thus $\mathcal{A}(s)$ is state-dependent and already accounts for stochastic energy via the safety multiplier $\alpha$, respecting charger availability and preventing policies from choosing actions that would strand the truck.

\end{document}
